% Draft chapter. Bibliography keys are placeholders pending integration with the
% thesis-wide .bib file. A source checklist appears at the end of this file.

\chapter{Introduction}
\label{chap:introduction}

\section{Problem and motivation}
\label{sec:introduction_problem}

Online misinformation changes as people repeat, reinterpret, and contest it.
This process matters when a scientific proposal becomes attached to an established conspiracy narrative.
Solar geoengineering provides such a case.
Public discussion of stratospheric aerosol intervention and Harvard's SCoPEx project has intersected with claims about chemtrails, weather control, and elite manipulation \parencite{debnath2023conspiracy}.
These claims can affect how citizens interpret research, institutions, and climate policy.
The relevant problem is therefore not only whether false content travels.
It is also how network structure and social identity shape factual change during transmission.

Online firestorm theory offers a social account of rapid, affective, and clustered communication.
Pfeffer, Zorbach, and Carley describe an online firestorm as a sudden volume of negative word of mouth directed at a person, company, or group \parencite{pfeffer2014firestorms}.
Their Outlook identifies seven interrelated factors: speed and volume, binary choices, network clusters, unrestrained information flow, lack of diversity, cross-media dynamics, and network-triggered decision processes.
The theory explains why particular communication environments can intensify collective reactions.
It does not specify a generative agent model or a quantitative tipping rule.

Geoengineering discourse supplies an empirical setting in which several of these mechanisms may matter.
Debnath and colleagues analysed 814,924 English-language tweets containing the geoengineering hashtag between 2009 and 2021 \parencite{debnath2023conspiracy}.
They identified a prominent chemtrails discourse alongside climate-action and environmental-concern themes.
They also documented temporal changes around solar-radiation-management projects and governance events.
Their study maps observed discourse at scale.
It does not provide controlled counterfactuals in which the same message enters different networks or identity compositions.

Agent simulation can provide such control, but it introduces a separate validity problem.
Classical cascade models represent diffusion through states or activation probabilities \parencite{kempe2003influence}.
They can isolate structural mechanisms, but they do not directly represent claim-level changes in natural language.
Large language model agents can transform text while retaining persona and interaction context.
Their outputs, however, reflect prompts, model behaviour, and scoring choices rather than observed human conduct.
Generative social simulations must therefore expose their assumptions and restrict claims to the evidence that their design supports \parencite{sargent2013verification}.

Maurya and colleagues provide the immediate technical basis for this thesis.
Their LASS workshop paper at ACM CIKM 2025 introduced an auditor--node framework for persona-conditioned rewrites on linear chains \parencite{maurya2025simulating}.
The framework scores factual degradation through a misinformation index (MI) and aggregates scores through a misinformation propagation rate (MPR).
That study used general-news domains, fixed chains of up to 30 rewrites, and a discrete question-answering auditor.
Its framework and published findings are prior work.
This thesis neither reintroduces those metrics nor reuses its crime and commercial results as new evidence.

\section{Research gap}
\label{sec:introduction_gap}

The three preceding research strands leave a specific methodological gap.
Pfeffer et al.\ provide a theory of firestorm conditions without a generative testbed.
Debnath et al.\ provide an observational account of geoengineering discourse without controlled interventions on network structure or identity composition.
Maurya et al.\ provide a generative misinformation instrument without a climate corpus, a graph-topology comparison, or an explicit mapping to firestorm theory.
The gap addressed here lies at their intersection.
It concerns how to design and audit a controlled climate-firestorm simulation without treating simulation output as observed social behaviour.

The proposal framed this intersection as six factors: valence, surprise, identity alignment, network clustering, information echo, and temporal acceleration.
Those labels are not Pfeffer et al.'s original seven-factor list.
This thesis treats them as a computational remapping.
It also records the parts that the simulation does not implement.
Cross-media dynamics has no media-broadcast layer.
Binary platform choices have no direct analogue because agents may forward, reinterpret, or drop a message.
Surprise remains a single-node drip seed, and the temporal horizon remains fixed.
The study therefore evaluates selected mechanisms rather than claiming to reproduce the complete firestorm theory.

Measurement creates a second gap inside the simulation.
The prior discrete auditor counts omitted and contradicted answers.
The present campaign also uses a continuous factual-recovery score.
These instruments share an approximate range but not a common measurement scale.
A difference between their outputs may reflect scoring design rather than a different social process.
The analysis must therefore keep discrete and continuous MI separate.

Empirical grounding creates a third constraint.
The public Debnath artefacts available to this study did not yield hydrated tweet text or conversation cascades.
The thesis consequently uses theory-faithful behavioural-profile reductions.
It also uses a 63-node hashtag co-occurrence graph as a structural fallback.
Neither object is an empirical user cluster or a retweet network.
This distinction prevents the simulation from being described as a digital twin of Twitter.

\section{Aim and analytical approach}
\label{sec:introduction_aim}

This thesis aims to develop and evaluate an exploratory artificial-society testbed for climate-misinformation propagation.
The testbed extends the prior auditor--node framework with climate and geoengineering seeds, Debnath-grounded behavioural profiles, and configurable network topologies.
It maps the firestorm theory to simulation controls, held conditions, and measured outcomes.
It then tests how identity composition, topology, and auditor design shape measured factual degradation.

The main campaign crosses eight graph topologies with two identity designs and two scoring modes.
The topologies are a linear chain, ring, Erdős--Rényi graph, small-world graph, scale-free graph, echo chamber, polarised graph, and hierarchy.
The homogeneous arm assigns one of 12 behavioural profiles to all eight nodes.
The heterogeneous arm assigns one of six eight-profile mixtures.
Each configuration processes six climate or geoengineering articles over eight ticks.
The resulting design contains 1,728 article-level cells.
All cells use \texttt{gpt-4o-mini}, one graph seed, and one stochastic run.

The analysis reports MI and mean node MPR within each scoring mode.
It defines \(k^{*}\) as the first tick at which mean MI exceeds 3 and remains above 3 for every later observed tick.
This quantity is a right-censored simulator clock.
It is not an estimate of ignition time on social media.
The analysis also retains 292 cells with at most one scored event as hatched dead cascades.
It excludes those cells from live-cell means but does not recode them as zero.

A separate D-net analysis places the same agents on the 63-node hashtag graph.
This analysis compares simulated cascade shape with the documented fallback structure.
It also compares homogeneous conspiracy profiles with a mixed profile assignment on the same edge set.
The comparison concerns graph structure and simulated auditor scores.
Debnath et al.\ report no empirical MI or MPR against which those scores could be validated.

\section{Research questions}
\label{sec:introduction_rqs}

The study addresses four research questions.
Each question follows from a design choice that the available evidence can evaluate.

\begin{enumerate}
    \item[\textbf{RQ1}] \textbf{Holding identity arm and topology labels fixed, how do the continuous and dual-discrete auditor instruments change estimated misinformation severity, \(k^{*}\) occurrence, and topology ordering?}

    RQ1 treats scoring mode as a measurement question.
    It compares modes without averaging their values or interpreting their gap as a third misinformation measure.

    \item[\textbf{RQ2}] \textbf{Within one scoring mode, does explicit conspiracy composition explain misinformation severity and \(k^{*}\) occurrence more consistently than the coarse homogeneous or heterogeneous label?}

    RQ2 compares persona families and named mixtures.
    It treats composition as substantive and H or He as an assignment description.

    \item[\textbf{RQ3}] \textbf{How does topology condition the same-mode relationship between persona composition, live-cell MI, and irreversible \(k^{*}\) within the eight-tick horizon?}

    RQ3 compares the eight generated topologies.
    It reports dead-cell rates with live-cell means because cascade survival changes the observed population.

    \item[\textbf{RQ4}] \textbf{What can the 63-node hashtag-based D-net comparison establish about structural similarity and identity-conditioned factual degradation, and where does that comparison fail?}

    RQ4 separates empirical graph properties from simulated content scores.
    It tests a documented fallback rather than claiming reconstruction of Debnath et al.'s tweet cascades.
\end{enumerate}

These questions support descriptive comparisons.
They do not support population estimates or causal effects.
The single stochastic run also precludes confirmatory inference about topology, identity, or scoring mode.

\section{Contributions}
\label{sec:introduction_contributions}

This thesis makes four bounded contributions.

\begin{enumerate}
    \item It provides an explicit mapping from Pfeffer et al.'s seven firestorm factors to the simulation's varied, held, measured, and absent elements.
    The mapping corrects the proposal's six-factor shorthand and retains cross-media dynamics as an unimplemented boundary.

    \item It constructs a documented climate and geoengineering simulation layer for the prior auditor--node framework.
    This layer contains six campaign articles, 12 homogeneous behavioural profiles, six heterogeneous mixtures, eight topologies, and reproducible configuration files.
    The profiles reduce themes reported by Debnath et al.; they are not clusters estimated from 814,924 tweets.

    \item It evaluates identity composition and network topology under two distinct factual-degradation instruments.
    The analysis preserves dead cascades, separates continuous from discrete scores, and distinguishes persona composition from the broad homogeneous--heterogeneous contrast.

    \item It links a 63-node hashtag co-occurrence structure to a custom simulation graph and audits the limits of that comparison.
    This analysis reports structural mismatch as well as similarity.
    It does not equate auditor MI with Twitter behaviour.
\end{enumerate}

These contributions concern operationalisation, experimental design, and transparent measurement.
They do not include the original auditor--node architecture, MI, MPR, or the published linear-chain findings.
They also do not establish a validated model of public response to geoengineering.

\section{Scope and limitations}
\label{sec:introduction_scope}

The empirical scope is deliberately narrow.
The campaign uses six authored or curated climate articles and a finite set of prompt-based behavioural profiles.
It runs eight-node graphs for eight ticks with one model and one graph seed.
The design changes text through LLM role-play.
It does not model stable human belief, individual psychology, or platform recommendation systems.

The data scope is narrower than the proposal anticipated.
The study did not hydrate the 814,924 tweet identifiers.
It did not estimate HDBSCAN profiles, reconstruct conversation trees, or reproduce Debnath et al.'s temporal series.
The hashtag graph records co-occurrence structure, not retweet transmission.
The simulation clock cannot be converted into minutes, days, or empirical hop distance.

The inferential scope follows these design limits.
The results describe one stochastic campaign.
They do not warrant significance claims, robust rankings, or general laws about homogeneous and heterogeneous groups.
The two auditors remain model-based measurements.
Their disagreement is evidence about the instrument, not an independent factual ground truth.
The \(k^{*}\) measure is conditional on the eight-tick observation window.
A value of ``none'' means that no irreversible crossing appeared within that window.

\section{Thesis roadmap}
\label{sec:introduction_roadmap}

Chapter~\ref{chap:related_work} develops the theoretical and empirical background.
It covers online firestorms, climate and geoengineering discourse, diffusion models, LLM-agent simulation, and automated factuality assessment.
Chapter~\ref{chap:theory} maps Pfeffer et al.'s seven factors to the thesis design.
Chapter~\ref{chap:methods} separates the inherited system from the climate-specific design and documents the articles, behavioural profiles, and hashtag graph.
Chapter~\ref{chap:methods} specifies the factorial campaign, both auditor modes, dead-cell rule, and \(k^{*}\).
Chapter~\ref{chap:results} answers RQ2 and RQ3 across the 1,728-cell campaign.
It then answers RQ4 through the D-net comparison.
Chapter~\ref{chap:discussion} interprets the findings against Pfeffer, Debnath, and the prior linear-chain study.
Chapter~\ref{chap:limitations} addresses validity, ethics, and reproducibility.
Chapter~\ref{chap:conclusion} answers the research questions and states the conditions required for stronger claims.

The thesis therefore proceeds from theory to operationalisation, from operationalisation to measurement, and from measurement to bounded interpretation.
Its central claim is methodological.
An LLM society can expose how climate-misinformation results depend on identity, topology, and the scoring instrument, but this campaign cannot stand in for the public it simulates.

% ---------------------------------------------------------------------------
% COMPANION SOURCE LIST FOR THIS DRAFT
% Add or verify these records in the thesis-wide bibliography before compiling.
%
% pfeffer2014firestorms
%   Pfeffer, J.; Zorbach, T.; Carley, K. M. (2014).
%   Understanding Online Firestorms: Negative Word-of-Mouth Dynamics in
%   Social Media Networks. Journal of Marketing Communications, 20(1--2),
%   117--128. DOI: 10.1080/13527266.2013.797778.
%   Supports: firestorm definition and original seven-factor Outlook.
%
% debnath2023conspiracy
%   Debnath, R. et al. (2023). Conspiracy Spillovers and Geoengineering.
%   iScience, 26(3), 106166. DOI: 10.1016/j.isci.2023.106166.
%   Supports: 814,924-tweet observational corpus, discourse themes, and
%   geoengineering context. It does not provide empirical MI or MPR.
%
% maurya2025simulating
%   Maurya, R. G. et al. (2025). Simulating Misinformation Propagation in
%   Social Networks using Large Language Models. 1st Workshop on LLM Agents
%   for Social Simulation (LASS), ACM CIKM 2025. arXiv:2511.10384.
%   Supports: prior auditor--node framework, MI/MPR, severity bands, and
%   30-hop linear-chain study. Cite as workshop work, not CIKM main track.
%
% kempe2003influence
%   Kempe, D.; Kleinberg, J.; Tardos, E. (2003). Maximizing the Spread of
%   Influence through a Social Network. KDD '03, 137--146.
%   DOI: 10.1145/956750.956769.
%   Supports: activation-based cascade-model contrast.
%
% tornberg2024validation
%   Törnberg, P. et al. (2024). Validation is the central challenge for
%   generative social simulation. Journal of Computational Social Science.
%   Verify final author list, volume, pages, and DOI before submission.
%   Supports: validation limits of generative social simulation.
%
% Internal evidence used for design statements, not bibliography entries:
%   thesisExperiment/PHASE2_PLAN.md
%   thesisExperiment/LOG.md
%   thesisExperiment/results_phase2/summary.json
%   thesisExperiment/analysis_phase2/SUMMARY.md
%   thesisExperiment/analysis_full/COMPARISONS.md
%   thesisExperiment/analysis_full/PFEFFER_DEBNATH_VS_SIM.md
%   thesisExperiment/analysis_phase2/PFEFFER_FIRESTORM.md
%   thesisExperiment/discovery/01_literature/_extract/proposal_pfeffer_v5.txt
%   thesisExperiment/discovery/01_literature/do_not_reprint.md
% ---------------------------------------------------------------------------
